\section{Related Work}

Standard augmentations preserve the original label, including region dropout
methods such as Cutout and Random Erasing~\citep{devries2017cutout,
zhong2020random} and policy-based methods such as AutoAugment, RandAugment,
and AugMix~\citep{cubuk2019autoaugment,
cubukRandAugmentPracticalAutomated2019,hendrycks2020augmix}. Pairwise
mixed-sample methods instead combine two examples and assign mixed targets:
Mixup~\citep{mixup} interpolates images and labels, while CutMix~\citep{cutmix}
uses spatial replacement with area-based labels. Later methods refine pairwise
composition or label estimation through resized patches, Fourier masks, grids,
saliency, class activation maps, attention, or token-level mixing~\citep{
qin2020resizemixmixingdatapreserved,harris2021fmixenhancingmixedsample,
baekGridMixStrongRegularization2021,kim2020puzzlemixexploitingsaliency,
huang2020snapmixsemanticallyproportionalmixing,
chen2021transmixattendmixvision,liu2023tokenmixrethinkingimagemixing}.
These methods improve label-preserving or two-source augmentation, whereas
TreemapMix constructs multi-source examples with a controllable distribution
over more than two visible labels.

\paragraph{Multi-image augmentation and label supervision:}
Multi-image augmentations compose more than two sources, including RICAP, which
patches four crops with area-weighted labels~\citep{
takahashiDataAugmentationUsing2020}, Mosaic augmentation in detection
pipelines~\citep{bochkovskiyYOLOv4OptimalSpeed2020}, Cut-thumbnail
~\citep{Xie_2021}, StackMix~\citep{chen2022stackmix}, and
DCutMix~\citep{jeong2022dcutmix}. Other methods generalize mixing beyond
spatial layouts, such as k-Mixup through optimal-transport interpolation
~\citep{greenewald2023kmixupregularizationdeeplearning} and MixMo through a
multi-input multi-output architecture~\citep{rame2021mixmo}. TreemapMix is
closest to spatial multi-image methods such as RICAP and Mosaic, but samples a
single coherent $K$-way area distribution from a Dirichlet distribution and
realizes it as an area-proportional treemap layout; unlike k-Mixup and MixMo,
it remains spatially explicit and uses a standard single-input model interface.
Most spatial mixing methods treat visible area as a soft-label
weight, although area may be an imperfect proxy for semantic evidence. Related
work studies confidence and target behavior under mixup-style training through
SmoothMixup~\citep{jeong2021smoothmixtrainingconfidencecalibratedsmoothed},
RegMixup~\citep{pinto2023regmixupmixupregularizersurprisingly}, and
RankMixup~\citep{noh2023rankmixuprankingbasedmixuptraining}. Plackett-Luce
ranking has also been used for classification in Plackett--Luce Distillation (PLD)~\citep{pld}, but in a
knowledge-distillation setting where a teacher model provides the class ranking.
In contrast, TreemapMix induces the ranking directly from the mixed input: the
Dirichlet-sampled region areas define the area-induced ordering over target
classes. TreemapMix-SCE follows the area-weighted probability interpretation,
while TreemapMix-PL converts the same areas into ordinal supervision with a
Plackett-Luce ranking objective.
